% ===========================================================================
%  第16章：自抗扰控制与在线辨识
% ===========================================================================

\section{自抗扰控制与在线参数辨识}
\label{sec:adrc_onlineid}

第\ref{sec:cascade_chapter}章的级联架构以PID内环+在线$\mathbf{B}_{\mathrm{true}}$分配
为核心控制路径。固件另提供两条辅助路径以增强鲁棒性：
\begin{enumerate}[label=(\arabic*)]
    \item \textbf{自抗扰控制（ADRC）}：以扩张状态观测器（ESO）实时估计并补偿
          总扰动（模型误差+外扰+交叉耦合），作为速率内环的PID替代方案；
    \item \textbf{在线参数辨识（OnlineID）}：以遗忘因子递推最小二乘（RLS）
          在线估计惯量比与总扰动，驱动自适应增益调度。
\end{enumerate}
本章给出两者的数学推导、离散化实现与稳定性约束，
与固件\texttt{TandemVec\_ADRC.h}和\texttt{TandemVec\_OnlineID.h}一一对应。

% ---------------------------------------------------------------------------
\subsection{二阶扩张状态观测器}
\label{sec:eso2}

\subsubsection{被控对象模型}

对单轴角速度动力学，将惯量不确定性、交叉耦合与外扰统一归入"总扰动"$f$：
\begin{equation}
    \dot{\omega} = b_0\,u + f
    \label{eq:adrc_plant}
\end{equation}
其中$u$为控制输入（角加速度指令，rad/s$^2$），
$b_0$为标称控制增益（默认$b_0=1$，即$u$与$\dot\omega$同量纲），
$f$为包含模型误差$\Delta b\cdot u$、交叉耦合与外扰的总扰动。
将$f$扩充为第三个状态$x_3=f$（假设$\dot f\approx0$），得到扩张状态空间：
\begin{equation}
    \begin{cases}
        \dot{x}_1 = x_2 + b_0\,u \\
        \dot{x}_2 = x_3 \\
        \dot{x}_3 \approx 0
    \end{cases}
    \qquad
    y = x_1 = \omega
    \label{eq:eso_plant}
\end{equation}

\subsubsection{ESO结构与带宽参数化}

采用Gao (2006)的带宽参数化形式，观测器增益由单一参数$\omega_o$（观测器带宽）决定：
\begin{equation}
    \boxed{
    \begin{aligned}
        \dot{\hat{x}}_1 &= \hat{x}_2 + b_0\,u + \beta_1(y - \hat{x}_1) \\
        \dot{\hat{x}}_2 &= \hat{x}_3 + \beta_2(y - \hat{x}_1) \\
        \dot{\hat{x}}_3 &= \beta_3(y - \hat{x}_1)
    \end{aligned}
    }
    \qquad
    \beta_1 = 3\omega_o,\quad
    \beta_2 = 3\omega_o^2,\quad
    \beta_3 = \omega_o^3
    \label{eq:eso_continuous}
\end{equation}
其中$\hat{x}_1\to\omega$（角速度估计），$\hat{x}_2\to\dot\omega$（角加速度估计），
$\hat{x}_3\to f$（总扰动估计）。
误差动力学特征多项式为$(s+\omega_o)^3$，观测误差以$\omega_o$为带宽指数收敛。

\subsubsection{固件简化：二阶ESO}

固件\texttt{TandemVec\_ADRC.h}采用二阶ESO（仅估计$\omega$和$f$，不单独估计$\dot\omega$），
适用于控制律仅需扰动补偿而不需角加速度估计的场景：
\begin{equation}
    \boxed{
    \begin{aligned}
        \dot{z}_1 &= z_2 - 2\omega_o\,e + b_0\,u \\
        \dot{z}_2 &= -\omega_o^2\,e
    \end{aligned}
    }
    \qquad e = z_1 - y
    \label{eq:eso2_continuous}
\end{equation}
其中$z_1\to\omega$，$z_2\to f$（总扰动）。
误差动力学特征多项式$(s+\omega_o)^2$，增益$\beta_1=2\omega_o$、$\beta_2=\omega_o^2$。

Euler离散化（控制周期$\Delta t$）：
\begin{equation}
    \begin{aligned}
        z_1 &\mathrel{+}= (z_2 - 2\omega_o\,e + b_0\,u)\,\Delta t \\
        z_2 &\mathrel{+}= (-\omega_o^2\,e)\,\Delta t
    \end{aligned}
    \label{eq:eso2_discrete}
\end{equation}

\textbf{实现要点}：ESO中的$u$必须取\textbf{上一拍实际输出}$u_{k-1}$而非当前指令$u_k$。
若用当前指令，ESO会将控制效果误判为扰动并加以对消，导致闭环失效。

\subsubsection{ADRC控制律}

基于ESO估计的PD控制+扰动对消：
\begin{equation}
    \boxed{
    u = \frac{1}{b_0}\left(\dot{\omega}_{\mathrm{ref}}
        + K_p(\omega_{\mathrm{ref}} - z_1) - z_2\right)
    }
    \label{eq:adrc_law}
\end{equation}
其中$K_p=\omega_c$（控制器带宽），$\dot{\omega}_{\mathrm{ref}}$为参考轨迹微分
（悬停定点时为零）。$-z_2/b_0$项实现对总扰动的实时对消。

闭环特征多项式$s+\omega_c$（一阶），带宽$\omega_c$决定跟踪速度。
推荐整定：$\omega_c=6$\,rad/s，$\omega_o=8$\,rad/s（$\omega_o/\omega_c\approx1.3$）。

% ---------------------------------------------------------------------------
\subsection{ADRC稳定性约束}
\label{sec:adrc_stability}

ESO带宽$\omega_o$受执行器带宽$\omega_{\mathrm{act}}\approx1/\tau_m$约束。
对本构型$\tau_m=0.28$\,s，$\omega_{\mathrm{act}}\approx3.6$\,rad/s。
仿真扫频结果：
\begin{itemize}
    \item $\omega_o\leq12$\,rad/s（$\approx3.3\times\omega_{\mathrm{act}}$）：
          全$\omega_c\in[2,20]$区间稳定；
    \item $\omega_o=50$\,rad/s：所有$\omega_c$均发散——
          ESO试图估计超出执行器响应能力的快变扰动，产生正反馈。
\end{itemize}
工程准则：$\omega_o\leq3\omega_{\mathrm{act}}$，$\omega_c\leq\omega_o/2$。
固件默认$\omega_o=8$、$\omega_c=6$满足此约束。

% ---------------------------------------------------------------------------
\subsection{遗忘因子递推最小二乘（RLS）}
\label{sec:rls}

\subsubsection{辨识模型}

对式(\ref{eq:adrc_plant})的离散化，以角加速度为观测量：
\begin{equation}
    y_k = \boldsymbol{\varphi}_k^T\,\boldsymbol{\theta} + \varepsilon_k
    \label{eq:rls_model}
\end{equation}
其中：
\begin{itemize}
    \item $y_k = \dot\omega_k$：由陀螺仪数值微分获得的实际角加速度（rad/s$^2$）；
    \item $\boldsymbol{\varphi}_k = [\alpha_{\mathrm{cmd},k},\;1]^T$：回归向量；
    \item $\boldsymbol{\theta} = [b,\;d]^T$：待估参数——
          $b$为实际惯量比（$I_{\mathrm{true}}/I_{\mathrm{nom}}$的倒数），
          $d$为常值扰动（含重心偏移、稳态耦合）；
    \item $\varepsilon_k$：测量噪声与未建模动态。
\end{itemize}

\subsubsection{RLS递推方程}

带遗忘因子$\lambda\in(0,1]$的标准RLS：
\begin{equation}
    \boxed{
    \begin{aligned}
        \hat{y}_k &= \boldsymbol{\varphi}_k^T\,\hat{\boldsymbol{\theta}}_{k-1} \\
        e_k &= y_k - \hat{y}_k \\
        \mathbf{S}_k &= \mathbf{P}_{k-1}\,\boldsymbol{\varphi}_k \\
        \mathbf{K}_k &= \frac{\mathbf{S}_k}{\lambda + \boldsymbol{\varphi}_k^T\,\mathbf{S}_k} \\
        \hat{\boldsymbol{\theta}}_k &= \hat{\boldsymbol{\theta}}_{k-1}
            + \mathbf{K}_k\,e_k \\
        \mathbf{P}_k &= \frac{1}{\lambda}
            \left(\mathbf{I} - \mathbf{K}_k\,\boldsymbol{\varphi}_k^T\right)
            \mathbf{P}_{k-1}
    \end{aligned}
    }
    \label{eq:rls_update}
\end{equation}
初始条件：$\hat{\boldsymbol{\theta}}_0=[1,0]^T$（标称惯量、零扰动），
$\mathbf{P}_0=100\,\mathbf{I}_2$（高初始不确定性）。
遗忘因子$\lambda=0.99$：等效窗口约$1/(1-\lambda)=100$步（200\,Hz下约0.5\,s），
使估计跟踪慢时变参数（如燃油消耗、载荷变化）而不过度响应噪声。

\subsubsection{激励充分性检测}

RLS在激励不足时协方差矩阵退化（"协方差爆炸"或"参数漂移"）。
固件采用指数移动均方值检测激励水平：
\begin{equation}
    \sigma_k^2 = 0.95\,\sigma_{k-1}^2 + 0.05\,\alpha_{\mathrm{cmd},k}^2
    \label{eq:excitation}
\end{equation}
当$\sigma_k^2\leq4.0$（即RMS$<2$\,rad/s$^2$）时跳过更新，
保持$\hat{\boldsymbol{\theta}}$和$\mathbf{P}$不变。
物理意义：悬停稳态下控制指令接近零，此时辨识无信息量，
仅在机动（阶跃、振荡、扰动）期间更新。

\subsubsection{参数更新限幅与安全防护}

为防止异常观测（传感器故障、碰撞）导致参数跳变：
\begin{equation}
    |\hat{b}_k - \hat{b}_{k-1}|
    \leq \max(|\hat{b}_{k-1}|,\;0.1)\times r_{\lim}
    \label{eq:rate_limit}
\end{equation}
其中$r_{\lim}=0.05$（每步最大变化5\%）。
硬约束：$\hat{b}\in[0.3,\;3.0]$，超出范围或检测到NaN时触发完全重置
（$\hat{\boldsymbol{\theta}}\leftarrow[1,0]^T$，$\mathbf{P}\leftarrow100\mathbf{I}$）。

% ---------------------------------------------------------------------------
\subsection{自适应增益调度}
\label{sec:adaptive_scheduling}

辨识得到的惯量比$\hat{b}$驱动速率内环增益调度：
\begin{equation}
    \boxed{
    K_p^{\mathrm{adapt}} = \frac{K_p^{\mathrm{nom}}}{\sqrt{\hat{b}}},
    \qquad
    \alpha_{\max}^{\mathrm{adapt}} = \alpha_{\max}^{\mathrm{nom}}\cdot\hat{b}
    }
    \label{eq:gain_scheduling}
\end{equation}
物理依据：内环闭环带宽$\omega_{\mathrm{BW}}=K_p/I$。
若真实惯量$I_{\mathrm{true}}=\hat{b}\cdot I_{\mathrm{nom}}$，
则$K_p/\sqrt{\hat{b}}$使带宽按$\sqrt{\hat{b}}$缩放——
惯量增大时降低增益避免过冲，惯量减小时提高增益维持响应速度。
$\alpha_{\max}$按$\hat{b}$正比缩放，确保力矩限幅与真实惯量匹配。

% ---------------------------------------------------------------------------
\subsection{重心偏移在线提取}
\label{sec:cg_extraction}

俯仰通道的常值扰动$d_{\mathrm{pitch}}$主要来源是重心（CG）偏离标称位置。
在近悬停工况下（油门30\%--70\%，三轴角速率$<15$\,deg/s），
俯仰积分项$I_{\mathrm{pitch}}$（rad/s）稳态时恰好补偿CG偏移产生的常值力矩：
\begin{equation}
    M_{\mathrm{CG}} = I_y\cdot K_i\cdot I_{\mathrm{pitch}}
    \label{eq:cg_moment}
\end{equation}
设总推力$T_{\mathrm{tot}}=2k_T\omega_0^2$，CG偏移量（mm）为：
\begin{equation}
    \Delta x_{\mathrm{CG}} = \frac{M_{\mathrm{CG}}}{T_{\mathrm{tot}}}\times1000
    \label{eq:cg_offset}
\end{equation}
固件以低通滤波（$\alpha=0.05$）平滑估计值，限幅$\pm20$\,mm，
供遥测监视和后续配平补偿使用（当前为观测模式，不自动修正控制律）。

% ---------------------------------------------------------------------------
\subsection{与级联架构的集成关系}
\label{sec:adrc_integration}

固件中ADRC与OnlineID的角色定位：
\begin{itemize}
    \item \textbf{主路径}（\texttt{flight\_control.cpp}）：
          \texttt{PositionPID}（位置式PID v3）姿态外环（纯P）+
          \texttt{PositionPID}速率内环（PI）+
          静态非线性控制分配（$\mathbf{M}=\mathbf{I}\boldsymbol{\alpha}$，
          \texttt{allocateMoments} BTRUE策略，上一拍工作点Jacobian
          = gain-scheduling，非NDI/INDI）；
    \item \textbf{ADRC}（\texttt{TandemVec\_ADRC.h}）：
          已完成实现与宿主机验证，\textbf{未接入生产控制链}。
          作为速率内环的PID替代方案储备，
          待需要常值扰动无积分对消能力时评估接入；
    \item \textbf{OnlineID}（\texttt{TandemVec\_OnlineID.h}）：
          已接入生产链但为\textbf{纯观测模式}——
          每拍调用\texttt{step()}估计$\hat{b}$、$\hat{d}$、$\Delta x_{\mathrm{CG}}$，
          结果写入\texttt{id\_*}全局遥测量供地面站与黑匣子记录。
          自适应增益建议$K_p/\sqrt{\hat{b}}$已计算（\texttt{id\_kp\_suggest}），
          但\textbf{未写回PID增益}，不影响任何控制输出。
          安全策略：第一步只观测不闭环，待多架次数据确认估计稳定后，
          再人工决定是否闭环接入（第二步）；
    \item \textbf{include/算法库}（\texttt{TandemVec\_CascadeCtrl.h}、
          \texttt{TandemVec\_AttitudeCtrl.h}、\texttt{TandemVec\_RateCtrl.h}）：
          平台无关（无Arduino依赖），用于\texttt{test\_host/}宿主机回归测试
          与算法验证，不是当前生产控制链的组成部分。
          生产链的姿态/速率环由\texttt{IncrementalPID}实例实现，
          两者在小角度下行为等价，但生产PID额外提供积分项与微分滤波。
\end{itemize}

% ---------------------------------------------------------------------------
\subsection{参考文献}

\begin{itemize}
    \item Han, J. (2009). \textit{From PID to Active Disturbance Rejection Control}.
          IEEE Trans. Industrial Electronics, 56(3), 900--906.
    \item Gao, Z. (2006). \textit{Scaling and Bandwidth-Parameterization based
          Controller Tuning}. Proc. ACC, 6--pp.
    \item Åström, K.J. \& Wittenmark, B. (1995). \textit{Adaptive Control}.
          2nd ed., Addison-Wesley.
    \item Ljung, L. (1999). \textit{System Identification: Theory for the User}.
          2nd ed., Prentice Hall.
\end{itemize}
